\subsection{容量完备、阈值分裂与账本--尺度障碍}\label{subsec:conclusion-capacity-threshold-shadow-ledger-scale-obstruction}
沿用定理 \ref{thm:pom-oracle-capacity-stieltjes-inversion-mellin}、定理 \ref{thm:xi-capacity-tail-histogram-moment-complete-statistic}、定理 \ref{thm:xi-time-part62df-capacity-curve-complete-inversion}、推论 \ref{cor:conclusion-renyi-universal-defect-oracle-rate-identification}、定理 \ref{thm:xi-full-microstate-exact-inversion-bitrate-threshold}、定理 \ref{thm:conclusion-frozen-phase-statistical-inversion-exponent-splitting}、定理 \ref{thm:conclusion-finite-prime-solenoid-terminal-object}、定理 \ref{thm:conclusion-finite-prime-localization-pontryagin-rigidity}、定理 \ref{thm:pom-resonance-mod2-finite-bit-automaticity} 与定理 \ref{thm:xi-resonance-closure-no-finite-lp-summability} 的记号。前述链条已经分别把容量曲线对纤维谱的完备编译、几乎全恢复与全精确反演的两类位率门槛、冻结相的统计/可逆指数分裂、局部化 prime-register 的一圆维 Lie 可见层，以及共振窗奇偶阴影的有限比特周期结构固定下来。把这些接口并置之后，还可进一步压出一条跨容量几何、峰值障碍、有限域阴影、prime-register 账本与多尺度复访成本的综合链。

\begin{theorem}[整数预算成功曲线的完备决定性]\label{thm:conclusion-integer-budget-curve-complete-profile}
\leanverified{paper\_conclusion\_integer\_budget\_curve\_complete\_profile}
固定分辨率 \(m\)，定义整数预算成功曲线
\[
\mathscr S_m(s)
:=
2^{-m}\sum_{x\in X_m}\min\bigl(d_m(x),s\bigr)
\qquad (s\in\ZZ_{\ge 0}).
\]
则整条函数 \(s\mapsto \mathscr S_m(s)\) 与下列对象彼此唯一决定：
\begin{enumerate}
  \item 纤维多重度直方图
  \[
  n_m(k):=\#\{x\in X_m:\ d_m(x)=k\};
  \]
  \item 连续容量曲线
  \[
  \mathcal C_m(u):=\sum_{x\in X_m}\min\bigl(d_m(x),u\bigr)\qquad (u\ge 0);
  \]
  \item 尾计数函数
  \[
  N_m(u):=\#\{x\in X_m:\ d_m(x)>u\}\qquad (u\ge 0);
  \]
  \item 全部正实阶幂和矩
  \[
  S_m(a):=\sum_{x\in X_m}d_m(x)^a\qquad (a>0).
  \]
特别地，若另一套折叠系统 \((\widetilde X_m,\widetilde d_m)\) 满足
\[
\mathscr S_m(s)=\widetilde{\mathscr S}_m(s)
\qquad (\forall\, s\in\ZZ_{\ge 0}),
\]
则两者在该分辨率上具有相同的纤维多重度多重集，并因而具有相同的连续容量、尾计数与全部正矩谱。
\end{enumerate}
\end{theorem}

\begin{proof}
由定义直接有
\[
\mathcal C_m(s)=2^m\mathscr S_m(s)
\qquad (s\in\ZZ_{\ge 0}).
\]
因此整数预算成功曲线与连续容量曲线在全部整数节点上的取值列等价。再由定理 \ref{thm:xi-time-part62df-capacity-curve-complete-inversion} 与定理 \ref{thm:xi-capacity-tail-histogram-moment-complete-statistic}，整数节点上的容量列唯一决定纤维直方图；反向方向显然，因为直方图直接给出 \(\mathscr S_m(s)\)。

进一步，由定理 \ref{thm:pom-oracle-capacity-stieltjes-inversion-mellin}，
\[
\mathcal C_m(u)=\int_0^u N_m(v)\,\dd v,
\qquad
S_m(a)=a\int_0^\infty u^{a-1}N_m(u)\,\dd u.
\]
故一旦容量曲线或直方图被确定，尾计数函数与全部正矩谱亦随之唯一确定；反之，直方图显然可由尾计数恢复。于是题述各类数据彼此唯一决定。证毕。
\end{proof}

\begin{corollary}[预算曲线的复杂度同型判据]\label{cor:conclusion-budget-curve-complexity-isomorphism}
\leanverified{paper\_conclusion\_budget\_curve\_complexity\_isomorphism}
若两套折叠系统在固定分辨率 \(m\) 上具有相同的整数预算成功曲线 \(s\mapsto \mathscr S_m(s)\)，则任一仅依赖纤维多重度谱的最优恢复量都在两者之间一致；特别地，由连续容量、尾计数或矩谱所导出的全部纤维复杂度剖面完全相同。
\end{corollary}

\begin{proof}
由定理 \ref{thm:conclusion-integer-budget-curve-complete-profile}，相同的整数预算成功曲线唯一确定同一纤维多重度多重集。任何只以该多重度谱为输入的恢复量因而必然一致。证毕。
\end{proof}

\begin{theorem}[容量曲线的整 Mellin 决定性]\label{thm:conclusion-capacity-entire-mellin-determinacy}
\leanverified{paper\_conclusion\_capacity\_entire\_mellin\_determinacy}
固定分辨率 \(m\)。设纤维大小的不同取值为
\[
1\le D_1<\cdots<D_r,
\qquad
n_i:=\#\{x\in X_m:\ d_m(x)=D_i\}.
\]
定义整 Mellin 生成函数
\[
\mathfrak Z_m(s)
:=
\sum_{x\in X_m}d_m(x)^s
=
\sum_{i=1}^{r} n_iD_i^s
\qquad (s\in\CC).
\]
则连续容量曲线
\[
\mathcal C_m(u):=\sum_{x\in X_m}\min\bigl(d_m(x),u\bigr)
\qquad (u\ge 0)
\]
唯一决定整个函数 \(\mathfrak Z_m\)。反之，\(\mathfrak Z_m\) 在任一具有聚点的子集上的取值唯一决定 \(\mathcal C_m\)。
\end{theorem}

\begin{proof}
由定理 \ref{thm:conclusion-integer-budget-curve-complete-profile}，\(\mathcal C_m\) 唯一决定纤维多重度直方图，因而唯一决定有限指数和
\[
\mathfrak Z_m(s)=\sum_{i=1}^{r} n_i e^{s\log D_i}.
\]

反向设两组纤维直方图给出的整函数 \(\mathfrak Z_m,\widetilde{\mathfrak Z}_m\) 在某个具有聚点的集合 \(E\subset\CC\) 上满足
\[
\mathfrak Z_m(s)=\widetilde{\mathfrak Z}_m(s)
\qquad (s\in E).
\]
则差函数 \(\mathfrak Z_m-\widetilde{\mathfrak Z}_m\) 是在 \(E\) 上有聚点零集的整函数，故由恒等定理知
\[
\mathfrak Z_m\equiv \widetilde{\mathfrak Z}_m
\qquad \text{于 }\CC.
\]
特别地，对一切 \(a>0\) 都有
\[
\sum_{x\in X_m}d_m(x)^a
=
\mathfrak Z_m(a)
=
\widetilde{\mathfrak Z}_m(a).
\]
再由定理 \ref{thm:xi-capacity-tail-histogram-moment-complete-statistic}，全部正实阶矩唯一决定纤维直方图，从而唯一决定 \(\mathcal C_m\)。证毕。
\end{proof}

\begin{corollary}[固定分辨率下的热解析性]\label{cor:conclusion-fixed-resolution-no-thermal-singularity}
\leanverified{paper\_conclusion\_fixed\_resolution\_no\_thermal\_singularity}
在定理 \ref{thm:conclusion-capacity-entire-mellin-determinacy} 的记号下，定义有限尺度配分函数与压力
\[
Z_m(\beta):=\mathfrak Z_m(-\beta)=\sum_{x\in X_m}d_m(x)^{-\beta},
\qquad
P_m(\beta):=\log Z_m(\beta).
\]
则 \(Z_m\) 是整函数，且对每个实数 \(\beta\) 都满足
\[
Z_m(\beta)>0.
\]
因此 \(P_m\) 在实轴上实解析。特别地，固定分辨率 \(m\) 上不存在任何实轴热奇点、真实相切割或有限尺度相变；全部此类现象至多只能出现在 \(m\to\infty\) 的极限过程中。
\end{corollary}

\begin{proof}
由定理 \ref{thm:conclusion-capacity-entire-mellin-determinacy}，
\[
Z_m(\beta)=\sum_{i=1}^{r} n_i e^{-\beta\log D_i}
\]
是有限指数和，故为整函数。若 \(\beta\in\RR\)，则每一项 \(n_i e^{-\beta\log D_i}\) 都严格为正，从而
\[
Z_m(\beta)>0.
\]
于是 \(P_m(\beta)=\log Z_m(\beta)\) 在 \(\RR\) 上可取单值实对数并成为实解析函数。证毕。
\end{proof}

\begin{theorem}[顶层纤维的 Mellin 剥离重构]\label{thm:conclusion-mellin-top-layer-peeling-reconstruction}
\leanverified{paper\_conclusion\_mellin\_top\_layer\_peeling\_reconstruction}
在定理 \ref{thm:conclusion-capacity-entire-mellin-determinacy} 的设定下，令
\[
\mathfrak Z_m^{(0)}:=\mathfrak Z_m,
\]
并对 \(0\le j\le r-1\) 递归定义
\[
D_{r-j}
:=
\lim_{\sigma\to+\infty}\bigl(\mathfrak Z_m^{(j)}(\sigma)\bigr)^{1/\sigma},
\qquad
n_{r-j}
:=
\lim_{\sigma\to+\infty}\mathfrak Z_m^{(j)}(\sigma)D_{r-j}^{-\sigma},
\]
\[
\mathfrak Z_m^{(j+1)}(s)
:=
\mathfrak Z_m^{(j)}(s)-n_{r-j}D_{r-j}^{\,s}.
\]
则该递归恰在 \(r\) 步后终止，并精确恢复全部纤维谱 \(\{(D_i,n_i)\}_{i=1}^{r}\)。因此，容量曲线经由 \(\mathfrak Z_m\) 不仅唯一确定直方图，而且允许以逐层 Laplace 剥离的显式方式恢复顶层纤维结构，而无需求助于 Hankel 行列式或 Prony 线性系统。
\end{theorem}

\begin{proof}
对 \(j=0\)，有
\[
\mathfrak Z_m^{(0)}(s)=\sum_{i=1}^{r}n_iD_i^s.
\]
由于 \(D_r\) 是最大底数，当 \(\sigma\to+\infty\) 时，
\[
\mathfrak Z_m^{(0)}(\sigma)
=
D_r^{\sigma}\bigl(n_r+o(1)\bigr).
\]
故
\[
\lim_{\sigma\to+\infty}\bigl(\mathfrak Z_m^{(0)}(\sigma)\bigr)^{1/\sigma}=D_r,
\qquad
\lim_{\sigma\to+\infty}\mathfrak Z_m^{(0)}(\sigma)D_r^{-\sigma}=n_r.
\]
减去该顶层项后得到
\[
\mathfrak Z_m^{(1)}(s)=\sum_{i=1}^{r-1}n_iD_i^s.
\]

若某一步已有
\[
\mathfrak Z_m^{(j)}(s)=\sum_{i=1}^{r-j}n_iD_i^s,
\]
则同理其最大底数为 \(D_{r-j}\)，从而
\[
\mathfrak Z_m^{(j)}(\sigma)
=
D_{r-j}^{\sigma}\bigl(n_{r-j}+o(1)\bigr)
\qquad (\sigma\to+\infty),
\]
于是题述两个极限分别恢复 \(D_{r-j}\) 与 \(n_{r-j}\)。减去该项后，剩余函数再次成为项数减少一的有限指数和。有限步归纳即得结论。证毕。
\end{proof}

\begin{theorem}[体量阈值与峰值阈值的分离]\label{thm:conclusion-volume-peak-threshold-separation}
\leanverified{paper\_conclusion\_volume\_peak\_threshold\_separation}
在二进折叠主线中，记
\[
\rho_{\mathrm{vol}}:=\log_2\frac{2}{\varphi},
\qquad
\rho_{\mathrm{peak}}:=\frac{\alpha_\ast}{\log 2},
\qquad
\alpha_\ast:=\lim_{m\to\infty}\frac1m\log M_m.
\]
则存在两条来源不同的必要位率门槛：
\begin{enumerate}
  \item 若整数预算序列 \(B_m\) 满足
  \[
  \mathrm{Succ}^{\mathrm{bin}}_m(B_m)\longrightarrow 1,
  \]
  则
  \[
  \liminf_{m\to\infty}\frac{B_m}{m}\ge \rho_{\mathrm{vol}}.
  \]
  \item 设 \(\Lambda_m:\Omega_m\to\{0,1\}^{B_m}\) 为任意侧信息标签。若对每个充分大的 \(m\) 与每个 \(x\in X_m\)，映射
  \[
  \omega\longmapsto \bigl(\Fold_m(\omega),\Lambda_m(\omega)\bigr)
  \]
  在纤维 \(\Fold_m^{-1}(x)\) 上都为单射，则
  \[
  2^{B_m}\ge M_m,
  \qquad
  \liminf_{m\to\infty}\frac{B_m}{m}\ge \rho_{\mathrm{peak}}.
  \]
\end{enumerate}
特别地，若 \(\alpha_\ast>\log(2/\varphi)\)，则
\[
\rho_{\mathrm{peak}}-\rho_{\mathrm{vol}}>0.
\]
因此全局体量障碍与局部峰值障碍一般并不重合。
\end{theorem}

\begin{proof}
第一条正是推论 \ref{cor:conclusion-renyi-universal-defect-oracle-rate-identification} 的位率下界。

对第二条，固定充分大的 \(m\)，取一条最大纤维 \(F_x:=\Fold_m^{-1}(x)\) 满足 \(|F_x|=M_m\)。若 \(\omega\mapsto(\Fold_m(\omega),\Lambda_m(\omega))\) 在 \(F_x\) 上单射，则不同的 \(\omega\in F_x\) 必对应不同的 \(B_m\)-比特标签；故标签空间 \(\{0,1\}^{B_m}\) 至少含有 \(M_m\) 个不同值，即
\[
2^{B_m}\ge M_m.
\]
取对数并除以 \(m\)，再令 \(m\to\infty\)，便得
\[
\liminf_{m\to\infty}\frac{B_m}{m}\ge \frac{\alpha_\ast}{\log 2}=\rho_{\mathrm{peak}}.
\]
最后，若 \(\alpha_\ast>\log(2/\varphi)\)，则
\[
\rho_{\mathrm{peak}}-\rho_{\mathrm{vol}}
=
\frac{\alpha_\ast-\log(2/\varphi)}{\log 2}>0.
\]
证毕。
\end{proof}

\begin{corollary}[冻结相中主导障碍的峰值定位]\label{cor:conclusion-frozen-phase-peak-obstruction-localization}
\leanverified{paper\_conclusion\_frozen\_phase\_peak\_obstruction\_localization}
若 \(\alpha_\ast^{(2)}<\alpha_\ast\) 且 \(a>a_{\mathrm c}\)，则冻结相中的全部平衡统计指数塌缩为 \(g_\ast\)，而精确可逆恢复的临界线性位率仍固定为 \(\alpha_\ast/\log 2\)。因而峰值阈值对应的并非可忽略的边缘异常层，而是冻结统计所集中到的主导纤维层。
\end{corollary}

\begin{proof}
定理 \ref{thm:conclusion-frozen-escort-shannon-renyi-collapse} 已给出：当 \(a>a_{\mathrm c}\) 时，escort 测度在最大纤维层 \(A_m^{\max}\) 上指数集中。
定理 \ref{thm:conclusion-frozen-phase-statistical-inversion-exponent-splitting} 已给出
\[
\mathfrak h_\theta(a)=g_\ast
\qquad (1\le \theta\le \infty),
\qquad
\beta_{\mathrm{inv}}(a)=\frac{\alpha_\ast}{\log 2}.
\]
因此冻结相把全部平衡统计指数压缩到 \(g_\ast\)，却不改变精确可逆恢复所需的峰值位率。换言之，最难恢复的顶端纤维同时也是冻结统计的主导承载层。证毕。
\end{proof}

\begin{proposition}[共振窗奇偶阴影的常数内存塌缩]\label{prop:conclusion-parity-shadow-constant-memory-collapse}
定义向量值奇偶阴影
\[
\mathbf s(m)
:=
\bigl(S_9(m)\bmod 2,\dots,S_{17}(m)\bmod 2\bigr)\in \FF_2^9.
\]
则存在某个 \(m_0\) 与函数
\[
F:\ZZ/16\ZZ\longrightarrow \FF_2^9
\]
使得对所有 \(m\ge m_0\) 都有
\[
\mathbf s(m)=F(m\bmod 16).
\]
因此，任一仅依赖 \(\mathbf s(m)\) 的流式审计器都可由一个内部状态数不超过 \(16\) 的确定有限自动机实现；特别地，这类奇偶阴影不可能充当容量曲线或纤维直方图的完备统计量。
\end{proposition}
\leanverified{paper\_conclusion\_parity\_shadow\_constant\_memory\_collapse}

\begin{proof}
由定理 \ref{thm:pom-resonance-mod2-finite-bit-automaticity}，对每个 \(q\in\{9,\dots,17\}\) 都存在
\[
k_q=\left\lceil\log_2 b_q\right\rceil
\]
使得 \(S_q(m)\bmod 2\) 在充分大 \(m\) 上只依赖于 \(m\bmod 2^{k_q}\)。表 \ref{tab:fold_collision_charpoly_mod2_shadow_q9_17} 给出
\[
b_q\in\{4,6,8,10\},
\]
故对全部 \(q=9,\dots,17\) 都有 \(k_q\le 4\)。于是九维向量 \(\mathbf s(m)\) 在充分大 \(m\) 上统一只依赖于 \(m\bmod 16\)，即存在所述映射 \(F\)。

因而只需记录当前余类 \(m\bmod 16\) 即可逐步更新 \(\mathbf s(m)\)，确定有限自动机的内部状态数至多为 \(16\)。另一方面，定理 \ref{thm:conclusion-integer-budget-curve-complete-profile} 表明容量曲线或纤维直方图在固定分辨率上是完备统计量，而 \(\mathbf s(m)\) 在长期段仅取有限多个值，故不可能充当此类完备统计量。证毕。
\end{proof}

\begin{proposition}[圆维轴与 prime-register 账本轴的严格正交]\label{prop:conclusion-circle-dimension-prime-ledger-strict-orthogonality}
对任意有限素数集 \(S\subset\mathbb P\)，记
\[
G_S:=\ZZ[S^{-1}],
\qquad
\Sigma_S:=\widehat{G_S}.
\]
则成立：
\leanverified{paper\_conclusion\_circle\_dimension\_prime\_ledger\_strict\_orthogonality}
\begin{enumerate}
  \item 存在短正合列
  \[
  0\longrightarrow \prod_{p\in S}\ZZ_p
  \longrightarrow \Sigma_S
  \longrightarrow \TT
  \longrightarrow 0,
  \qquad
  \cdim^\ast(G_S)=1.
  \]
  \item \(\Sigma_S\cong \Sigma_T\) 当且仅当 \(S=T\)。
  \item 不存在函数 \(F:[0,\infty)\to\NN\) 使得对所有有限 \(S\subset\mathbb P\) 都有
  \[
  |S|\le F\bigl(\cdim^\ast(G_S)\bigr).
  \]
\end{enumerate}
换言之，连续圆维只记录唯一的一圆相位轴，而 prime-register 的规模与支撑完全滞留在全不连通账本核中。
\end{proposition}

\begin{proof}
第一条由定理 \ref{thm:conclusion-finite-prime-solenoid-terminal-object} 与定理 \ref{thm:cdim-star-z1s-dual-extension} 直接给出。第二条正是定理 \ref{thm:conclusion-finite-prime-localization-pontryagin-rigidity} 的内容。

若第三条中的函数 \(F\) 存在，则由 \(\cdim^\ast(G_S)=1\) 可得
\[
|S|\le F(1)
\]
对一切有限 \(S\subset\mathbb P\) 成立。但有限素数集的基数可以任意大，矛盾。故不存在此类函数。

最后，由第二条可知真正决定 \(S\) 的正是核
\[
\prod_{p\in S}\ZZ_p,
\]
而不是恒等于 \(1\) 的圆维。证毕。
\end{proof}

\begin{corollary}[相位预算与 prime-ledger 预算不可互偿]\label{cor:conclusion-phase-budget-and-prime-ledger-budget-noncompensable}
设 \(I\) 是定义在有限素数局部化族 \(\{G_S=\ZZ[S^{-1}]\}\) 上的一个不变量，并满足：
\begin{enumerate}
  \item 可加性；
  \item 有限扩张不计费；
  \item 有限指数不变。
\end{enumerate}
则 \(I\) 在该族上只能测到 circle-dimension 方向，因而必经由常值
\[
\cdim^\ast(G_S)=1
\]
因子化；特别地，\(I\) 不可能恢复 \(S\) 的 prime-ledger 支撑。于是，任何 exact implementation 若既要记录连通相位预算，又要记录素数账本成本，就必须额外并列一个残差预算轴，二者之间不存在互偿原则。
\end{corollary}
\leanverified{paper\_conclusion\_phase\_budget\_and\_prime\_ledger\_budget\_noncompensable}
\begin{proof}
由命题 \ref{prop:conclusion-circle-dimension-prime-ledger-strict-orthogonality}，
\[
\cdim^\ast(G_S)=1
\qquad
(\forall\,S\subset\mathbb P\ \text{有限}),
\]
而真正决定 \(S\) 的信息却完整滞留在全不连通核
\[
\prod_{p\in S}\ZZ_p
\]
之中。另一方面，正文关于 circle-dimension 公理唯一性的结果表明：凡满足可加性、有限扩张不计费与有限指数不变的此类不变量，都只能经由 circle-dimension 因子化。因此这类 \(I\) 在 \(\{G_S\}\) 上至多是常值，不能恢复 prime-ledger 支撑。

若仍试图把相位预算与素数账本预算合并为单一可加预算，则该预算必须同时既经由 circle-dimension 因子化、又能区分不同的 \(S\)，与上段矛盾。故相位预算与 prime-ledger 残差预算只能并列存在，而不能彼此抵偿。证毕。
\end{proof}

\begin{theorem}[固定二维 \texorpdfstring{$2$}{2}-进 ambient 的充分性与有限秩同态账本的绝对不足]\label{thm:conclusion-prime-register-fixed-2adic-ambient-vs-finite-ledger}
\leanverified{paper\_conclusion\_prime\_register\_fixed\_2adic\_ambient\_vs\_finite\_ledger}
设 \(H_{\mathrm{can}}\) 为 canonical 有限时深 prime-register 子单子。则存在忠实单子表示
\[
\Psi:H_{\mathrm{can}}\hookrightarrow \operatorname{Aff}\!\bigl(T_2(E_{\min})\bigr),
\qquad
T_2(E_{\min})\cong \ZZ_2^2,
\]
亦即固定二维 \(2\)-进 Tate 模块已足以承载全部有限深度 canonical prime-register 动力学。

另一方面，不存在任何有限生成交换可消去加法账本 \(L\) 使得 cofinal 的乘法结构 \(\NN^\times\) 能忠实嵌入
\[
(\ZZ\oplus L,+)
\qquad\text{或}\qquad
(\NN^k,+)
\]
对任何有限 \(k\) 成立。因而 prime-multiplicative 机制的真正障碍不在 ambient 空间维数，而在要求把乘法语义保持为有限秩加法账本同态。
\end{theorem}
\begin{proof}
主文 H.160 已构造出 canonical 仿射作用
\[
\Psi(r,k)(x)=B^k x+\operatorname{ev}_B(r)
\]
并证明其在 \(H_{\mathrm{can}}\) 上忠实，其中底层空间正是
\[
T_2(E_{\min})\cong \ZZ_2^2.
\]
故固定二维 \(2\)-进 Tate 模块确已足够。

另一方面，主文 H.130--H.131 已证明
\[
\NN^\times\not\hookrightarrow (\NN^k,+)
\qquad(\forall\,k<\infty),
\]
且更强地，对任意有限生成交换群 \(L\) 都不存在忠实幺半群嵌入
\[
\NN^\times\hookrightarrow (\ZZ\oplus L,+).
\]
因此，有限维 ambient realization 的存在，与有限秩 additive ledger realization 的存在是两件根本不同的事；前者成立，后者失败。证毕。
\end{proof}

\begin{theorem}[一寄存器可序列化与线性同态债务的同时成立]\label{thm:conclusion-prime-register-one-register-serialization-vs-linear-hom-debt}
\leanverified{paper\_conclusion\_prime\_register\_one\_register\_serialization\_vs\_linear\_hom\_debt}
记
\[
S_r:=\{p_1,\dots,p_r\},
\qquad
M_r:=\NN^\times_{S_r}.
\]
则有精确维数差公式
\[
\operatorname{cdim}^{+}_{\mathrm{hom}}(M_r)=\frac r2,
\qquad
\operatorname{cdim}_{\mathrm{impl}}(M_r)=\frac12,
\qquad
\operatorname{cdim}^{+}_{\mathrm{hom}}(M_r)-\operatorname{cdim}_{\mathrm{impl}}(M_r)=\frac{r-1}{2}.
\]
换言之，同一对象族可以在实现层被无损序列化到单寄存器，而在结构层却需要线性增长的加法寄存器数才能保持幺半群语义。
\end{theorem}
\begin{proof}
主文已证明：任意单射幺半群同态
\[
\Phi_r:M_r\hookrightarrow (\NN^{k_r},+)
\]
必满足 \(k_r\ge r\)，并且该下界可达；因此
\[
\operatorname{cdim}^{+}_{\mathrm{hom}}(M_r)=\frac r2.
\]

另一方面，只要 \(r\ge 1\)，\(M_r\) 就是可数无限集。主文关于实现层半圆维的定义与其可数无限对象分类给出
\[
\operatorname{cdim}_{\mathrm{impl}}(M_r)=\frac12.
\]
更具体地，命题 13.4059 说明任意有限描述对象族都可通过可计算单射
\[
\mathrm{enc}:A^\ast\hookrightarrow \NN
\]
无损序列化到单寄存器；但该编码并不保持幺半群运算，因此不会降低结构层的同态圆维。

两式相减即得
\[
\operatorname{cdim}^{+}_{\mathrm{hom}}(M_r)-\operatorname{cdim}_{\mathrm{impl}}(M_r)=\frac{r-1}{2}.
\]
证毕。
\end{proof}


\begin{proposition}[共振闭包预算的 \texorpdfstring{$\ell^2$}{l2} 不可压缩性]\label{prop:conclusion-resonance-budget-no-l2-envelope}
\leanverified{paper\_conclusion\_resonance\_budget\_no\_l2\_envelope}
记 \(C_\varphi(u)\) 为黄金共振闭包常数序列。则不存在常数 \(K>0\) 与序列 \(w=(w_u)_{u\ge 1}\in \ell^2(\NN)\) 使得
\[
C_\varphi(u)\le K\,w_u
\qquad (\forall\,u\ge 1).
\]
因此，任何试图以平方可和尾预算统一吸收全部多尺度共振复访成本的方案都不可能成立。
\end{proposition}

\begin{proof}
若存在此类 \(K\) 与 \(w\in \ell^2(\NN)\)，则
\[
\sum_{u\ge 1}C_\varphi(u)^2
\le
K^2\sum_{u\ge 1}w_u^2
<
\infty.
\]
这与定理 \ref{thm:xi-resonance-closure-no-finite-lp-summability} 在 \(p=2\) 的情形矛盾。故所述平方可和包络不存在。证毕。
\end{proof}

\begin{theorem}[全整数预算成功曲线的完全充分性]\label{thm:conclusion-budget-curve-exact-tail-difference-reconstruction}
\leanverified{paper_conclusion_budget_curve_exact_tail_difference_reconstruction_seeds}
固定分辨率 \(m\)。定义全整数预算成功曲线
\[
\mathrm{Succ}_m(s):=
2^{-m}\sum_{x\in X_m}\min\bigl(d_m(x),s\bigr),
\qquad
s=0,1,\dots,2^m,
\]
并记
\[
C_m(s):=2^m\mathrm{Succ}_m(s),
\qquad
\Delta_m(s):=C_m(s)-C_m(s-1)
\qquad (1\le s\le 2^m),
\]
以及纤维多重度直方图
\[
M_m(k):=\#\{x\in X_m:\ d_m(x)=k\}.
\]
则有精确公式
\[
\Delta_m(s)=\#\{x\in X_m:\ d_m(x)\ge s\},
\]
\[
M_m(k)=\Delta_m(k)-\Delta_m(k+1).
\]
因此整条成功曲线 \(s\mapsto \mathrm{Succ}_m(s)\) 唯一决定纤维多重度多重集 \(\{d_m(x):x\in X_m\}\)。
\end{theorem}

\begin{proof}
对任意固定整数 \(d\ge 0\) 与 \(s\ge 1\)，有恒等式
\[
\min(d,s)-\min(d,s-1)=\mathbf 1_{\{d\ge s\}}.
\]
将其应用到每条纤维多重度 \(d_m(x)\) 并对 \(x\in X_m\) 求和，得到
\[
\Delta_m(s)
=
\sum_{x\in X_m}
\Bigl(\min(d_m(x),s)-\min(d_m(x),s-1)\Bigr)
=
\#\{x\in X_m:\ d_m(x)\ge s\}.
\]
再用尾计数差分恒等式
\[
\#\{d_m\ge k\}-\#\{d_m\ge k+1\}
=
\#\{d_m=k\}
\]
即得
\[
M_m(k)=\Delta_m(k)-\Delta_m(k+1).
\]
于是 \(\mathrm{Succ}_m\) 唯一决定全部 \(M_m(k)\)，进而唯一决定纤维多重度多重集。证毕。
\end{proof}

\begin{corollary}[预算层即 Mellin--Rényi 因子化层]\label{cor:conclusion-budget-curve-mellin-renyi-factorization-layer}
\leanverified{paper_conclusion_budget_curve_mellin_renyi_factorization_layer}
设 \((\mathrm{Succ}_m)_{m\ge 1}\) 为全部分辨率上的整数预算成功曲线塔。则对任何仅依赖纤维多重度多重集的观测量，它都是完全充分统计量。特别地，下列对象都经由预算曲线塔唯一因子化：
\[
S_q(m)=\sum_{k\ge 1}k^q M_m(k)
\qquad (q>0),
\]
\[
C_m^{\mathrm{cont}}(T)=\sum_{k\ge 1}\min(k,T)M_m(k),
\]
\[
N_m(t)=\sum_{k\ge \lceil t\rceil}M_m(k).
\]
因而全部 Mellin 矩、连续容量曲线、尾谱、以及由这些对象诱导的压力谱 \(P_q\)、冻结阈值、基态简并指数 \(g_\ast\) 与容量指数包络，都经由预算曲线塔唯一决定。
\end{corollary}

\begin{proof}
由定理 \ref{thm:conclusion-budget-curve-exact-tail-difference-reconstruction}，预算曲线塔唯一决定全部直方图 \(M_m(k)\)。而题中各式都是对 \(M_m(k)\) 的显式函数，因此均由预算曲线塔唯一决定。项目文稿中压力、冻结阈值、基态简并指数与容量指数包络的定义只依赖于这些由直方图导出的量，故结论成立。证毕。
\end{proof}

\begin{proposition}[可逆容量指数的普适折角]\label{prop:conclusion-capacity-exponent-universal-fold-angle}
记
\[
\Gamma(\beta):=
\limsup_{m\to\infty}\frac1m\log C_m(\lfloor \beta m\rfloor),
\qquad
\beta\ge 0,
\]
并定义
\[
M_m^{\max}:=\max_{x\in X_m}d_m(x),
\qquad
\alpha_\ast:=\lim_{m\to\infty}\frac1m\log M_m^{\max},
\qquad
g_\ast:=\lim_{m\to\infty}\frac1m\log K_m.
\]
则项目文稿给出的普适下支满足
\[
\Gamma(\beta)\ge g_\ast+\min\{\alpha_\ast,\beta\log 2\}.
\]
若记
\[
\beta_c:=\frac{\alpha_\ast}{\log 2},
\]
则该下支
\[
\underline\Gamma(\beta):=g_\ast+\min\{\alpha_\ast,\beta\log 2\}
\]
在 \(\beta_c\) 处具有不可消除的折角：
\[
\underline\Gamma(\beta)=
\begin{cases}
g_\ast+\beta\log 2,& \beta\le \beta_c,\\[1mm]
g_\ast+\alpha_\ast,& \beta\ge \beta_c.
\end{cases}
\]
其左斜率为 \(\log 2\)，右斜率为 \(0\)。
\end{proposition}

\begin{proof}
这是项目文稿中容量指数双侧夹逼之下支的直接改写。由
\[
\beta\log 2=\alpha_\ast
\]
的唯一交点恰为
\[
\beta_c=\frac{\alpha_\ast}{\log 2},
\]
故 \(\underline\Gamma\) 的显式分段式立即成立，左右导数分别为 \(\log 2\) 与 \(0\)。证毕。
\end{proof}
\leanverified{paper\_conclusion\_capacity\_exponent\_universal\_fold\_angle}

\begin{theorem}[有界实现切片上的 bulk--boundary 有限完备审计]\label{thm:conclusion-bulk-boundary-joint-finite-complete-audit}
固定分辨率 \(m\)，并假设对象属于定理 \ref{thm:conclusion-toeplitz-endpoint-atom-dyadic-joint-verification} 的有界实现--端点谱隙切片。定义审计对象
\[
\mathcal A_{m,D,\varepsilon}
:=
\Bigl(
(\mathrm{Succ}_m(s))_{s=0}^{2^m},
\ (T_{2^k})_{k=0}^{K(D,\varepsilon)},
\ a_{2^{K(D,\varepsilon)}}
\Bigr).
\]
则 \(\mathcal A_{m,D,\varepsilon}\) 同时完成下列三件互异任务：
\begin{enumerate}
  \item \textbf{bulk 重构：}唯一恢复纤维多重度直方图 \(M_m(k)\)；
  \item \textbf{global feasibility：}判定全部 Toeplitz--PSD 约束；
  \item \textbf{boundary extraction：}以误差不超过 \(\varepsilon\) 的精度恢复端点原子质量 \(\mu(\{-1\})\)。
\end{enumerate}
\end{theorem}

\begin{proof}
由定理 \ref{thm:conclusion-budget-curve-exact-tail-difference-reconstruction}，第一分量 \((\mathrm{Succ}_m(s))_{s=0}^{2^m}\) 唯一恢复全部直方图 \(M_m(k)\)，从而完成 bulk 重构。由定理 \ref{thm:conclusion-toeplitz-endpoint-atom-dyadic-joint-verification}，第二分量 \((T_{2^k})_{k=0}^{K(D,\varepsilon)}\) 判定全层级 Toeplitz--PSD，可完成 global feasibility；第三分量 \(a_{2^{K(D,\varepsilon)}}\) 以误差不超过 \(\varepsilon\) 恢复 \(\mu(\{-1\})\)，完成 boundary extraction。三者合并即得结论。证毕。
\end{proof}

\begin{theorem}[dyadic 边界全息与端点原子抽取的双通道不可约分裂]\label{thm:conclusion-bulk-boundary-endpoint-dual-channel-irreducibility}
固定分辨率 \(m\)，并仍处于定理 \ref{thm:conclusion-bulk-boundary-joint-finite-complete-audit} 的有界实现--端点谱隙切片。则 faithful 的边界审计至少必须同时保留两条互补通道：
\[
H_m(U)
\qquad\text{与}\qquad
\mu(\{-1\}).
\]
更精确地，二元对
\[
\bigl(H_m(U),\mu(\{-1\})\bigr)
\]
在该切片上 jointly faithful；而任一只经由
\[
H_m(U)
\qquad\text{或}\qquad
\mu(\{-1\})
\]
单独因子化的边界审计都不是 faithful。
\end{theorem}
\begin{proof}
由定理 \ref{thm:conclusion-dyadic-boundary-code-dual-realization-parameters}，边界 Gödel 整数 \(H_m(U)\) 已足以唯一恢复 fixed-resolution 的 dyadic polycube，因而完整承载 bulk 边界几何。另一方面，定理 \ref{thm:conclusion-toeplitz-endpoint-atom-dyadic-joint-verification} 表明端点原子质量 \(\mu(\{-1\})\) 由独立的端点热核探针通道抽取，而定理 \ref{thm:conclusion-three-end-certificate-joint-sufficiency-minimality} 与命题 \ref{prop:conclusion-three-end-certificate-orthogonality} 又排除了用其余证书轴替代该线性代数端通道的可能。

因此，只保留 \(H_m(U)\) 会丢失端点原子信息；只保留 \(\mu(\{-1\})\) 则会丢失 dyadic bulk 的全息重构。两条通道合并时，前者负责 arithmetic holography，后者负责 endpoint extraction，故二元对 \(\bigl(H_m(U),\mu(\{-1\})\bigr)\) 在该切片上 jointly faithful。于是 faithful 的边界审计至少必须保留这两条互补通道。证毕。
\end{proof}

\leanverified{paper\_conclusion\_fiber\_tail\_spectrum\_renyi\_pressure\_upper\_bound}
\begin{theorem}[纤维尾谱的 R\'enyi--压力上界]\label{thm:conclusion-fiber-tail-spectrum-renyi-pressure-upper-bound}
定义尾计数函数
\[
N_m(\theta):=\#\{x\in X_m:\ d_m(x)\ge \theta\},
\]
其中 \(d_m(x)=|\Fold_m^{-1}(x)|\)。则对任意整数 \(q\ge 1\) 与任意 \(\gamma\ge 0\)，都有
\[
\limsup_{m\to\infty}\frac1m\log N_m(e^{\gamma m})
\le
\log r_q-q\gamma,
\]
其中
\[
r_q:=\lim_{m\to\infty}S_q(m)^{1/m},
\qquad
S_q(m):=\sum_{x\in X_m}d_m(x)^q.
\]
因此
\[
\limsup_{m\to\infty}\frac1m\log N_m(e^{\gamma m})
\le
\inf_{q\in \NN,\ q\ge 1}\bigl(\log r_q-q\gamma\bigr).
\]
\end{theorem}
\begin{proof}
由 Markov 型估计，
\[
N_m(e^{\gamma m})\,e^{q\gamma m}
\le
\sum_{x:\,d_m(x)\ge e^{\gamma m}}d_m(x)^q
\le
\sum_{x\in X_m}d_m(x)^q
=
S_q(m).
\]
两边取对数并除以 \(m\)，再令 \(m\to\infty\)，得到
\[
\limsup_{m\to\infty}\frac1m\log N_m(e^{\gamma m})
\le
\lim_{m\to\infty}\frac1m\log S_q(m)-q\gamma
=
\log r_q-q\gamma.
\]
对所有整数 \(q\ge 1\) 取下确界，即得第二式。证毕。
\end{proof}

\begin{conjecture}[纤维尾谱的大偏差精确公式]\label{conj:conclusion-fiber-tail-spectrum-legendre-sharp}
设文稿中的实参数 Perron 分支
\[
\Lambda(t):=\log \lambda(t+1)-\log |A_{\mathrm{out}}|
\qquad (t\ge 0)
\]
存在、可微且本质严格凸。则对每个处于暴露斜率内部的 \(\gamma\)，应有精确极限
\[
\lim_{m\to\infty}\frac1m\log N_m(e^{\gamma m})
=
\inf_{t\ge 0}\bigl(\log|A_{\mathrm{out}}|+\Lambda(t)-(t+1)\gamma\bigr).
\]
换言之，定理 \ref{thm:conclusion-fiber-tail-spectrum-renyi-pressure-upper-bound} 的上界在热力学良性区间内应当是 sharp 的；纤维尾谱的指数函数恰为实参数压力的 Legendre--Fenchel 对偶。

该猜想与文稿中的端点定理
\[
\lim_{q\to\infty}\frac{1}{q}\log r_q=\frac12\log\varphi
\]
一致；因此其右端的有效支撑应在
\[
\gamma\le \frac12\log\varphi
\]
处终止。
\end{conjecture}


\paragraph{统一读法}
上述综合链可压缩为
\[
\text{整数预算曲线}
\Longrightarrow
\text{纤维谱完备性}
\Longrightarrow
\text{体量/峰值阈值分离}
\Longrightarrow
\text{二维 ambient 充分而有限秩同态账本不足}
\Longrightarrow
\text{单寄存器实现与线性同态债务分裂}
\Longrightarrow
\text{奇偶阴影的有限状态塌缩}
\Longrightarrow
\text{圆维与 prime-register 正交}
\Longrightarrow
\text{bulk--boundary 联合有限完备审计}
\Longrightarrow
\text{共振复访预算的 }\ell^2\text{ 不可压缩性}.
\]
第一环给出复杂度的可测坐标，而且其一阶差分已经精确恢复尾计数与多重度直方图，因此全部 Mellin 矩、连续容量曲线、尾谱、压力谱与冻结阈值都经由预算层唯一因子化；第二环把全局体量障碍与局部峰值障碍分开，并进一步显影出容量指数下支在 \(\beta_c=\alpha_\ast/\log 2\) 处的普适折角；第三、第四环进一步说明：prime-register 的困难并不在于缺少一个固定低维 ambient，而在于任何保持乘法语义的有限秩加法账本都必然失败，因此实现层的一寄存器可序列化与结构层的线性同态债务严格并存；第五环说明低特征奇偶阴影只能提供有限状态投影；第六环排除了用单一连通圆维替代离散账本支撑的可能；第七环则把预算直方图、有限多个 dyadic Toeplitz 视界与单个端点探针拼接成 bulk--boundary 的联合有限完备审计；最后一环表明多尺度共振成本不能被平方可和尾项吸收。因而在本文固定的折叠口径下，逆向困难并不是单一复杂度标签，而是由容量几何、峰值纤维、ambient 实现/账本债务裂缝、有限域阴影、prime-register 核、bulk--boundary 审计耦合与尺度复访共同锁定的复合障碍。

\endinput
